../cictikz/src/cictikz/data/tex/cictikz_preamble.tex

% Sampling and aliasing: the same spectrum before and after a sample and hold.
%
% Two identical frequency axes, so the eye can compare them directly. The green
% hump is the wanted signal and survives; the four tones outside the Nyquist
% band fold inside it, which is the whole point of the slide.
%
% The redraw makes the folding arithmetically right, which the hand-drawn
% original only suggests. A tone at 1.2 fs lands on 0.2 fs and one at 0.7 fs
% lands on -0.3 fs, so the colours let you trace which tone went where. The
% relative order the original shows -- magenta ending up outside red -- is
% preserved, because that is what the arithmetic gives.
%
% No ckt_lib here: nothing in this figure is a circuit element.

\definecolor{axisblue}{rgb}{0.16,0.16,0.55}
\definecolor{specgreen}{rgb}{0.00,0.45,0.20}
\definecolor{tonered}{rgb}{0.85,0.22,0.10}
\definecolor{tonemag}{rgb}{0.90,0.20,0.75}

\newcommand{\half}{1.2}     % f_s/2
\newcommand{\xmax}{3.5}     % axis half length
\newcommand{\vh}{1.9}       % vertical axis height
\newcommand{\th}{1.05}      % tone height

% One frequency axis with its Nyquist ticks. #1 is the baseline.
\newcommand{\specaxis}[1]{%
  \draw[axisblue, thick, <->] (-\xmax,#1) -- (\xmax,#1);
  \draw[axisblue, thick, ->] (0,#1) -- (0,#1+\vh);
  \foreach \p/\lbl in {-2/{-f_{s}}, -1/{-f_{s}/2}, 1/{f_{s}/2}, 2/{f_{s}}} {
    \draw[axisblue, thick] (\p*\half,#1+0.13) -- (\p*\half,#1-0.13);
    \node[axisblue, anchor=north] at (\p*\half,#1-0.22) {$\lbl$};
  }
}

% The wanted signal, centred on DC. #1 is the baseline.
\newcommand{\bump}[1]{%
  \draw[specgreen, very thick] plot[smooth, tension=0.7] coordinates
    {(-1.05,#1) (-0.75,#1+0.80) (-0.35,#1+1.30) (0,#1+1.45)
     (0.35,#1+1.30) (0.75,#1+0.80) (1.05,#1)};
}

% A tone: colour, position, baseline.
\newcommand{\tone}[3]{\draw[#1, very thick, ->] (#2,#3) -- (#2,#3+\th);}

\begin{circuitikz}[american, transform shape]

  \def\ytop{3.9}
  \def\ybot{0}

  % ---- Before sampling: two pairs of tones outside the Nyquist band ----
  \specaxis{\ytop}
  \bump{\ytop}
  \tone{tonered}{-2.88}{\ytop}     % -1.2 fs
  \tone{tonemag}{-1.68}{\ytop}     % -0.7 fs
  \tone{tonemag}{1.68}{\ytop}
  \tone{tonered}{2.88}{\ytop}
  \node[anchor=west] at (1.45,\ytop+1.62) {Before sampling};

  % ---- After sampling: both pairs have folded inside ----
  \specaxis{\ybot}
  \bump{\ybot}
  \tone{tonered}{-0.48}{\ybot}     % 1.2 fs - fs = 0.2 fs
  \tone{tonemag}{-0.72}{\ybot}     % 0.7 fs - fs = -0.3 fs
  \tone{tonemag}{0.72}{\ybot}
  \tone{tonered}{0.48}{\ybot}
  \node[anchor=west] at (1.45,\ybot+1.62) {After sampling};

  % ---- The sample and hold doing it ----
  \def\bx{4.9}
  \def\by{2.0}
  \draw[axisblue, thick] (\bx,\by-0.5) rectangle ++(1.5,1.0);
  \node[axisblue] at (\bx+0.75,\by) {S/H};
  \draw[axisblue, thick, ->] (\bx-1.0,\by) -- (\bx,\by);
  \draw[axisblue, thick, ->] (\bx+1.5,\by) -- (\bx+2.5,\by);

\end{circuitikz}
\end{document}
